\section{Referee-controlled status and conditional geometric inputs}\label{sec:status}

The term \emph{proved} below is reserved for consequences for which the
stated input hypotheses and the displayed argument suffice.  An exact
rational verifier checks a finite predicate; it does not certify a
topological or analytic identification.  The companion independent
\texttt{REFEREE\_REPORT.md} (at the repository root) explains why the
following bridges must not be suppressed in an unconditional statement.
The substantial calculations of the original manuscript are preserved
in later sections as \emph{conditional derivations and proof attempts};
their mathematical existence in the archive is not independent review.

\begin{definition}[Unresolved geometric bridges]\label{def:bridges}
For a selected period $p$, we distinguish the following hypotheses
from the exact polynomial and rational inequalities.
\begin{description}
\item[(H1) Renormalization/nesting.]
The central $3p$-return has the asserted Chebyshev filled interval
with the two named endpoints, and any hypothetical competing
period-two renormalization of the \emph{central quadratic-like map}
is nested so that McMullen's almost-disjointness theorem can be
transported to its two little Julia sets; both inverse critical
points in the primeness contradiction belong to the proper
degree-two return domain.  No regularity of $\partial V_p$
stronger than local connectedness is silently assumed.

\item[(H2) Physical KKL ideal marking.]
On the specific two slit complements, the dynamically defined
KKL Hubbard tree, the hyperbolic prime-end/ideal boundary
marking, and the complete proper-path and subordination relation
are identified with the stated physically indexed slit-side
accesses.  In particular the inclusion's pulled marking
is compatible with the source covering marking and with
the refinement required by KKL Proposition~A.12.
This is a \emph{continuum class-identification} condition,
not merely a ten/twenty integer table.

\item[(H3) Full-path conformal comparison.]
The proposed embedded parent Schwarz--Christoffel
quadrilateral foliations consist of admissible simple
KKL A.1 Waves with their \emph{minimal} prescribed
bases and crests.  Every admissible source child Wave
has a lift into the proposed entire four-electrode
all-path envelope; the square-sum periodization is
admissible for \emph{all} such leaves, and the finite
source-excluded boxes remove positive energy without
changing their lengths.  The target lifted-path
envelope similarly contains \emph{every} marked parent
leaf.  The exact angle/AGM checks alone do not verify
these global inclusions.

\item[(H4) Peripheral-cover comparison.]
The universal-cover ideal interval lifts for the
four-hole KKL canonical WAD have the stated
disjoint-translate property, and the physical
Joukowski annuli compare to the corresponding
peripheral cyclic covers with the claimed transverse
width convention.  The finite ellipses themselves
are checked algebraically, but their interpretation
as bounds on \emph{all} ideal isotopy classes
uses this analytic input.
\end{description}
\end{definition}

These hypotheses are \emph{not} inferred from passing
Python/Lean checks or from prior in-repository audit
claims.  The named classical results provide plausible
tools to discharge them, but the referee has not verified
their complete applications in the manuscript.
Accordingly:
\begin{itemize}
\item The local root and Rouch\'e statements retain
their explicit, piecewise scope and written arithmetic/analytic
argument; their implementations are independently
\emph{executable} rather than machine proofs of the tail.
\item The complete physical marking conclusions are
asserted \emph{only under (H1)--(H2)}; the full
unbuffered canonical sandwich and strict
independently optimized defect are asserted
\emph{only under (H1)--(H3)}.  The reciprocal arithmetic
still gives strict $1/100$ once those inputs hold.
\item The buffered exclusion is conditional on
(H1), (H2), and (H4); no inner class is decided.
Any depth word's identification with a named
straightening suffix needs an independent \emph{full}
kneading argument beyond a table of real signs.
\end{itemize}

The conditional statements are useful to localize the exact
remaining proof tasks, but must \emph{not} be advertised
as unconditional new theorems or as KKL amplification.
The original first-version assertions remain unmodified
under \texttt{papers/tip\_family\_preprint\_20261116/}.
